\section{Selmer rigidity of a modular type-II point}
\label{sec:kato-rigidity}

Let \(K/\mathbb Q\) be the imaginary quadratic field of the window,
\(\theta\) its character, \(G=G_{\mathbb Q}\), and
\(\epsilon=\epsilon_\ell\) the cyclotomic character.
We prove the characteristic-zero Selmer vanishing used at a modular
type-II point. The arithmetic inputs are the finiteness theorem of
Kato and the adjoint vanishing theorem of Newton and Thorne. The
identification of the point as modular is a separate hypothesis.

\subsection{The local conditions and the two-dimensional summands}

For a continuous representation \(V\), write \(H^1_f(F,V)\) for
the classes of continuous cocycles whose restrictions to inertia
are coboundaries at every finite place away from \(\ell\), and whose
restrictions at \(v\mid\ell\) become coboundaries in
\(V\otimes_{\mathbb Q_\ell} B_{\rm cris}\). The latter condition says
that \(c(s)\otimes1=sx-x\) on the local absolute Galois group for
one element \(x\) of this tensor product with diagonal action.
Restrictions, finite direct sums and continuous equivariant linear
maps preserve these conditions. They agree with the conditions called
\(\mathrm{IsBK}\) and \(\mathrm{finiteCocycles}\) in the original
release. Choosing an embedding into a local algebraic closure selects a
decomposition group. Conjugating the embedding gives the same condition:
transport the local witness by the conjugating element and adjust by
the value of the cocycle on that element.

Let \(e\ge1\), \(k=2e+1\), and let
\(W:G\longrightarrow {\rm GL}_2(\overline{\mathbb Q}_\ell)\)
be continuous. Fix
\(\iota:\overline{\mathbb Q}_\ell\simeq\mathbb C\).
Suppose a normalized cuspidal newform \(f\) of weight \(k\), level
\(M_f\) and nebentypus \(\omega_f\) is attached to \(W(-e)\).
Here attachment requires semisimplicity, and for all but finitely
many primes \(p\), unramifiedness and the characteristic polynomial
\[
 X^2-a_p(f)X+\omega_f(p)p^{k-1}
 \quad\hbox{at geometric Frobenius, after applying }\iota.
\]
The original newform predicate imposes the good-prime eigenrelations,
normalization and exclusion of lower-level cusp forms with the same
good eigenvalues. Classical newform decomposition yields a primitive
eigenform with these good eigenvalues: decompose into old and new
eigensystems, choose a primitive constituent, and observe that a
constituent of smaller conductor would itself be a forbidden
lower-level form. At its primitive conductor the newvector line has
dimension one and the normalization fixes the scalar.

\paragraph{Proposition.}
Under these hypotheses
\[
 H^1_f(\mathbb Q,W)=H^1_f(\mathbb Q,W\otimes\theta)=0 .
 \tag{K1}
\]

\paragraph{Proof.}
Put \(E=F_\lambda\), where \(F\) is the coefficient field of \(f\)
and \(\lambda\) is the place induced by \(\iota^{-1}\).
Kato's representation \(V_E(f)\) has the displayed polynomial
at the inverse of arithmetic Frobenius (Example 13.3 of his paper).
Its semisimplicity and the Chebotarev density theorem give equality of
the traces of \(V_E(f)\) and \(W(-e)\) on all of \(G\) after scalar
extension. The characteristic-zero Brauer--Nesbitt theorem therefore
gives
\[
 W\simeq U\otimes_E\overline{\mathbb Q}_\ell,\qquad U=V_E(f)(e).
\]
In particular this is the positive Tate twist of Kato's cohomological
representation. The characteristic polynomial at geometric Frobenius
fixes the sign.

Take a \(G\)-stable \(\mathcal O_E\)-lattice \(T\) of \(U\).
Kato, \emph{\(p\)-adic Hodge theory and values of zeta functions of
modular forms}, Ast\'erisque 295 (2004), \S\S13.3, 14.1 and
Theorem 14.2(1), proves that for a finite abelian extension
\(L/\mathbb Q\) the discrete Selmer group of a stable lattice of
\(V_E(f)(r)\) is finite whenever \(1\le r\le k-1\) and \(2r\ne k\).
The discrete group is contained in \(H^1(L,U/T)\) and its local
conditions are the images of \(H^1_f(L_v,U)\).
Here \(r=e\) satisfies these inequalities and \(L=K\) is abelian.
The theorem applies to CM newforms as well.

We give the passage from discrete finiteness to vector-space
vanishing. The reduction map sends \(H^1_f(K,U)\) into this finite
group. Its image is divisible because its domain is an \(E\)-vector
space, so its image is zero. Let \(c\) represent any of its classes
and let \(\varpi\) be a uniformizer. For each \(j\ge1\), the reduction
of \(\varpi^{-j}c\) is therefore a coboundary. Lift the element of
\(U/T\) giving this coboundary to \(b_j\in U\). Then
\[
 c(s)-(s-1)\varpi^j b_j\in\varpi^jT\qquad(s\in G_K).
\]
The coboundaries form the image of the finite-dimensional space \(U\)
in the normed space of continuous functions \(G_K\to U\), with the
uniform norm from \(T\). This image is finite dimensional over the
complete local field \(E\), hence closed. The display gives uniform
convergence to \(c\), proving \(H^1_f(K,U)=0\).

We record the scalar-extension argument because the cocycle in
the statement takes values in an algebraic closure. Set
\(B=\overline{\mathbb Q}_\ell\). Continuous \(E\)-linear maps
\(\lambda:B\to E\) separate points. Indeed, on the line generated
by \(a\ne0\) prescribe \(\lambda(a)=1\), with norm bound
\(C=|a|^{-1}\). To extend a functional from a subspace \(D\)
of \(B\) to \(D+Ey\), choose its value at \(y\) in the balls
\[
 \{b\in E: |b-\lambda(x)|\le C|y-x|\},\qquad x\in D.
\]
The ultrametric inequality makes these balls nested. Their
intersection is nonempty: choose one of them, use its compactness
in the local field \(E\), and use the finite intersection property
for the smaller closed balls. The chosen value preserves the
bound on \(D+Ey\). The union of a chain of extensions preserves
the bound, so Zorn's lemma gives the required map on \(B\).

Apply every such \(\lambda\) coordinatewise to a finite cocycle
\(c:G_K\to U\otimes_E B\). This gives a continuous \(U\)-valued
cocycle satisfying the same local conditions, since the coefficient
map carries each local witness to a witness and extends by tensoring
with the identity on \(B_{\rm cris}\). Thus every projected cocycle
is a coboundary. To see that \(c\) is one, write the equations
\((U(s)-1)x=c(s)\) in an \(E\)-basis. Their left sides are rows
in the finite-dimensional \(E\)-dual of \(U\). Choose among all
rows a basis for their span, of dimension \(m\). Every other row is an \(E\)-linear
combination of the chosen ones. The same linear relation on the
right sides holds after every \(\lambda\), since the projected
equations have a solution. Separation of points gives each relation
in \(B\). The independent rows define a surjection to \(B^m\), so we can solve the basis equations.
This solution satisfies all equations. We have proved
\(H^1_f(K,W)=0\) without a restriction on the field generated by
the values of a cocycle.

Finally \(W\) and \(W\otimes\theta\) restrict to the same
representation on \(G_K\). Restriction preserves the local conditions.
Inflation--restriction identifies the kernel of restriction in
degree one with \(H^1({\rm Gal}(K/\mathbb Q),V^{G_K})\), for either
\(V\). Averaging over this group of order two kills that cohomology
in characteristic zero. Both restriction maps are injective,
proving (K1). \(\square\)

\subsection{The polarized tangent module}

We spell out the input supplied by the type-II classification.
Relative to the order \(0,1,2\), the matrix part of each \(r(s)\) is
\[
 P_s=\begin{pmatrix}a&0&b\\0&1&0\\c&0&d\end{pmatrix},\qquad
 M_s=\begin{pmatrix}a&b\\c&d\end{pmatrix}\in{\rm SL}_2 .
\]
Its component is \(\theta(s)\). At the chosen complex conjugation
\(c_\infty\), \(P_{c_\infty}=1\) and \(\theta(c_\infty)=-1\).
On \(G_K\) the attached two-dimensional representation is \(M_s\);
on the other coset it is \(M_sD\), with \(D={\rm diag}(-1,1)\).
The irreducibility supplied by the window is that of
\(W|_{G_{K(\mu_{\ell^\infty})}}\) over \(\overline{\mathbb Q}_\ell\).

The tangent module is the full space of three-by-three matrices.
On \(G_K\) the action is ordinary conjugation by \(P_s\). The chosen
complex conjugation acts by \(X\mapsto-J_3X^tJ_3\), where \(J_d\) is
the antidiagonal identity matrix of size \(d\). For \(t=1,-1\) set
\[
 p_t(X)=\binom{X_{01}+tX_{12}}{X_{21}-tX_{10}},\qquad
 s_t(u,v)=\begin{pmatrix}0&u&0\\-tv&0&tu\\0&v&0\end{pmatrix}.
\]
Let \(A(X)\) extract the outer corner on \(0,2\) and let \(i_A\)
insert it there. Write \(i_C(X_{11})\) for the central entry alone.
Direct substitution gives
\[
 X=i_A A(X)+\tfrac12s_+(p_+(X))
                +\tfrac12 s_-(p_-(X))+i_C(X_{11}).
 \tag{K2}
\]
Here \(+\) and \(-\) denote \(1\) and \(-1\). To verify the actions,
the two columns \((X_{01},X_{21})^t\) and
\((X_{12},-X_{10})^t\) both transform by \(M_s\) on \(G_K\).
For the second column use
\((X_{10},X_{12})\mapsto(X_{10},X_{12})M_s^{-1}\) and
\(\det M_s=1\). Complex conjugation exchanges these two columns with
the factor \(D\): the first becomes \(D\) times the second and the
second becomes \(D\) times the first. Thus their sum realizes \(W\),
their difference \(W\theta\). The center realizes \(\theta\).

On the outer corner \(G_K\) acts by conjugation. At \(c_\infty\)
the formula on two-by-two matrices is \(Y\mapsto-J_2Y^tJ_2\).
This is multiplication by \(-1\) on scalars. For
\(Y=\left(\begin{smallmatrix}u&v\\w&-u\end{smallmatrix}\right)\)
it equals \(DYD^{-1}\). Since \(2\) is invertible, the corner
decomposes as the trace-zero part with action \({\rm ad}^0 W\)
and a second copy of \(\theta\). Hence
\[
 {\rm ad}\,r \simeq
 {\rm ad}^0 W\ \oplus\ W\ \oplus\ (W\theta)\
                  \oplus\ \theta\ \oplus\ \theta .
 \tag{K3}
\]

\paragraph{Proposition.}
A modular type-II point with the block and irreducibility properties
above satisfies \(H^1_f(\mathbb Q,{\rm ad}\,r)=0\).

\paragraph{Proof.}
The newform \(f\) supplies a regular algebraic cuspidal
\({\rm GL}_2/\mathbb Q\) representation since \(k\ge2\).
In J.\ Newton and J.\ A.\ Thorne, \emph{Adjoint Selmer groups of
automorphic Galois representations of unitary type},
J.\ Eur.\ Math.\ Soc.\ 25 (2023), Theorem 5.4, the full adjoint
Selmer group of this representation vanishes if the form is non-CM,
or if its CM field is not contained in \(\mathbb Q(\mu_{\ell^\infty})\).
In the remaining possibility, with CM field \(L\) contained in that
cyclotomic field, the two-dimensional Galois representation is
induced from a character of \(G_L\). It splits into two characters
over \(G_L\) and hence is reducible on
\(G_{K(\mu_{\ell^\infty})}\subset G_L\). Twisting by \(\epsilon^e\)
preserves this reducibility, contrary to the specified hypothesis on
\(W\). The theorem therefore applies in every allowed case. Scalar
twists preserve the adjoint. The dual representation gives an
isomorphic adjoint via transposition, so either standard normalization
of an attached modular representation gives the same conclusion.
The trace-zero piece is a direct summand and
\(H^1_f(\mathbb Q,{\rm ad}^0W)=0\).

For the remaining character, restriction to \(K\) again injects,
and \(\theta|_{G_K}=1\). Over \(K\), the finite local condition
for the trivial representation \(E\) is unramified everywhere.
At places \(v\mid\ell\) this uses the standard local Bloch--Kato
fundamental exact sequence: since
\(D_{\rm dR}(E)/{\rm Fil}^0=0\), it identifies the finite local condition
with the cokernel of \(1-\varphi\) on \(D_{\rm cris}(E)\).
This group has dimension one over \(E\), since its kernel has
dimension one by unramified Frobenius descent. The unramified
line lies in it. Indeed an extension of the trivial representation
by itself with unramified action is crystalline (equivalently use
an unramified period solving \(\varphi^{f_v}(b)-b=1\), where
\(f_v\) is the residue degree).
A continuous additive character of \(G_K\) unramified at every
finite prime factors through the finite ideal class group by global
class field theory. Its image in \(E\) is zero. The separation
argument above also applies to \(B\)-valued classes. Thus
\(H^1_f(\mathbb Q,\theta)=0\). Apply (K1) and (K3). \(\square\)

\subsection{The Euler-system argument in the original release}

The archive reaches (K1) by constructing zeta classes. We describe
this separate derivation and the intermediate results it uses.
The original \(\mathrm{KatoTwoW}\) statement assumes an odd prime
\(\ell\) with \(\ell-1\nmid 2e\) and
\(\ell\nmid |\operatorname{disc}K|h_K\), an even
\(N>0\) prime to \(\operatorname{disc}K\), and a finite character
\(\psi\) factoring through \(\mathbb Q(\mu_N)\), of order prime
to \(\ell\) and quadratic on inertia above \(\ell\).
It imposes \(\Psi(c_\infty)=1\) for \(\Psi=\epsilon^e\psi\).
The level is \(N|\operatorname{disc}K|\) with ordered weights
\((e,-e)\), and the ordinary inertia polynomial after a finite
cyclotomic extension has roots \(\epsilon^e,\epsilon^{-e}\).
One assumes \(\det W=\theta\), inertia trace
\(\Psi+\Psi^{-1}\), and everywhere the residual trace is
\(\overline{\theta\Psi+\Psi^{-1}}\). At one place above \(\ell\)
there is a decomposition-group flag with diagonal inertia characters
\(\epsilon^e\chi,\epsilon^{-e}\chi\), with \(\chi\) factoring through
\(\mathbb Q(\mu_\ell)\), and \(\theta\) is trivial on this inertia
group. The irreducibility condition is the one in (K3).
The type-II point supplies these arithmetic properties through
the result \(\mathrm{used\_of\_alt}\). The archived version also
takes a rational cuspidal automorphic witness and the newform
attached to \(W(-e)\).

The comparison of the Galois row with quadratic base change and
its unramified local parameters provides a spherical compatibility
statement for the representation and its finite cyclotomic twists.
Comparison with an ordinary newform gives a newform \(g\) of
weight \(2e+1\), prime-to-\(\ell\) level, and
\(\lvert\iota^{-1}a_\ell(g)\rvert_\ell=1\). It is attached to
\[
 W\otimes\epsilon^{-e}(\det W)^{-1}\chi'
\]
for a character \(\chi'\) factoring through \(\mathbb Q(\mu_\ell)\).
For the selected flag with lower diagonal character \(\kappa_1\)
one has \(\kappa_1\epsilon^e(\det W)^{-1}\chi'=1\) on inertia.
The comparison also supplies nonexceptionality: for each admissible
decomposition-group flag, its first character differs from \(\epsilon\)
on some element of that decomposition group. These are the
unramified-source conditions in \(\mathrm{SrcL}\). After conjugation
to an ordinary integral basis they feed the construction
\(\mathrm{Final.zetaUr}\).

In this basis let \(A\) be the chosen place above \(\ell\) and \(D_A\)
its decomposition group. Write \({\rm Gr}\) for unramifiedness away
from \(\ell\) together with the condition that the lower coordinate
of the cocycle is a coboundary on \(D_A\). Write \({\rm Str}\) for
the same away condition together with being a coboundary in both
coordinates on \(D_A\). The Tate dual is
\(W'=W\otimes\epsilon(\det W)^{-1}\), via the alternating
form \(\omega(v,w)=v_0w_1-v_1w_0\). The zeta construction gives
an integral cocycle \(y\) unramified away from \(\ell\) whose lower
coordinate is not a coboundary on \(D_A\).

More precisely its norm relations supply the following data. There
is a finite exceptional set such that for \(n\ge1\), every auxiliary
prime \(q\) outside that set with \(\ell^n\mid q-1\), and every
surjective \(\psi_q:G\to\mathbb Z/\ell^n\) factoring through
\(\mathbb Q(\mu_q)\), one has integral vectors \(Y(s,x)\) for
\(x\in\mathbb Z/\ell^n\), a scalar \(P\), a unit \(\alpha\) of
absolute value one, and an integral vector \(b\), with
\begin{align*}
 Y(st,x)&=Y(s,x)+W(s)Y(t,x-\psi_q(s)),\\
 \sum_xY(s,x)&=Py(s)+(W(s)-1)b,\\
 P&\equiv \alpha\det(W({\rm Frob}_q)-1)\pmod{\ell^n}.
\end{align*}
The first family is continuous. The last congruence holds for each
admissible Frobenius lift, and the source imposes integral
inertia-coboundary witnesses for the family at all other primes
away from \(\ell\). Smoothed zeta classes, their integrality and
Hodge--Tate law, a nonvanishing Hecke projection, and a local
projection formula are the inputs to these assertions.

Here is the height contradiction that uses the norm relations.
Suppose that a strict dual cocycle \(z\) represents a nonzero class.
After enlarging the integral coefficient ring and multiplying \(y,z\)
by nonzero powers of \(\ell\), their values and their required local
witnesses are integral. Put
\[
 H=\ker\epsilon\cap\ker\det W,\qquad B_0=\ker\epsilon\cap\ker W.
\]
The cyclotomic-kernel lemmas in the release assert that \(y|_{B_0}=0\)
would give the ordinary condition on \(y\), and \(z|_{B_0}=0\)
would make \(z\) a global coboundary. The pairing and commutator
lemmas for the irreducible ordinary representation therefore supply
\(g,h\in B_0\) and \(\phi\in[H,H]\) such that
\[
 d'=\omega(y(h),z(g))-\omega(y(g),z(h))\ne0,\qquad
 d=\det(W(\phi)-1)\ne0.
\]
The elements are fixed for every precision \(n\).

For each \(n\ge1\), the cup product
\(\omega(y(s),W'(s)z(t))\) is locally a coboundary modulo
\(\ell^n\), by strictness at \(\ell\) and the unramified local cup
product away from \(\ell\). The global cochain gluing theorem
supplies an integral cochain \(u\) whose coboundary is this cup
product modulo \(\ell^n\), with the unramified-outside-\(S\)
condition. Define
\[
 T(s)=\det(W(s)-1)u(s)
  -\omega\bigl(\operatorname{adj}(W(s)-1)y(s),z(s)\bigr).
\]
The tame-lift theorem gives
\(\det(W(s)-1)T(s)\equiv0\bmod\ell^n\) on \([H,H]\).
Write \(c=[g,h]\). For clarity the commutator-height formula can
be seen directly. On \(B_0\) both actions are trivial, so \(y,z\)
restrict to homomorphisms and \(y(c)=z(c)=0\).
The cochain equation becomes
\[
 u(st)\equiv u(s)+u(t)-\omega(y(s),z(t))\pmod{\ell^n}
 \qquad(s,t\in B_0).
\]
Apply it to \(gh\) and \(hg\). Since \(gh=c\,hg\), it gives
\(u(c)\equiv d'\). The equation on \((\phi,c)\) now gives
\(u(\phi c)-u(\phi)\equiv d'\), because \(\epsilon(\phi)=1\).
The matrix and both cocycles at \(\phi c\) equal their values at
\(\phi\), hence \(T(\phi c)-T(\phi)\equiv dd'\). Both \(\phi\) and
\(\phi c\) belong to \([H,H]\), so the tame-lift theorem and
integrality of \(d\) give \(d^2d'\equiv0\bmod\ell^n\).
This holds for every \(n\) and contradicts \(d^2d'\ne0\).
Thus the strict dual group vanishes. The arithmetic input to
this calculation consists of the kernel and pairing lemmas,
local-to-global cochain gluing and the tame-lift height congruence.

The release upgrades strict dual vanishing to ordinary dual
vanishing by reciprocity for integral cup products and local
duality. After subtracting a coboundary, the lower coordinate
of an ordinary dual class is zero on \(D_A\). The local orthogonality
alternative makes either the lower coordinate of \(y\) a coboundary
or this adjusted dual class strict. The first option contradicts the
zeta class, and strict dual vanishing settles the second. Global
Selmer duality then gives ordinary vanishing for \(W\): a nonzero
ordinary class for \(W\) would produce dual finite-level classes
of unbounded order, and the compactness argument for these classes
produces a nonzero ordinary dual class.

Finally a finite class for \(W\) is ordinary. Its lower projection
is crystalline. For an upper triangular inertia matrix with diagonal
entries \(a,b\), the normal form gives \(a^m=\epsilon^{em}\)
for some \(m>0\) and \((ab)^2=1\). Thus
\((b\epsilon^e)^{2m}=1\) on inertia. The lower character has
weight \(-e\) in the cyclotomic-exponent convention. The
one-dimensional crystalline finite-class vanishing in negative weight
gives a coboundary on \(D_A\), as required. Applying this construction
to \(W\) and \(W\det W\) gives precisely (K1): the determinant
is quadratic and factors through a finite cyclotomic field, so the
spherical compatibility statement over all twists also applies to
the latter representation.

The original declarations for this route include
\(\mathrm{katoTwoW\_of\_printsH}\), \(\mathrm{twoS}\),
\(\mathrm{vanishOrdL\_of\_zetaUr}\), \(\mathrm{es\_vanish}\),
and \(\mathrm{StrictKill.main}\). Its type-II tangent proof projects
a cocycle by (K2), applies polarized-adjoint rigidity to the corner
and the nonpositive weight result to the center, then reassembles
the four coboundaries. The two propositions above supply a
classical derivation of the entire vanishing from the stated point
hypotheses. Establishing those hypotheses for all points of the
deformation ring, including production of \(f\), belongs to the
automorphy and classification arguments.
